\documentclass[12pt, twoside]{report}

%% comment this line for handouts. Uncomment for lecture notes
%\newcommand{\lectureNoteMode}{}

\raggedbottom

\usepackage{color, xifthen, amsthm, amssymb, amsmath, tikz, pgfplots}
\pgfplotsset{compat=1.18}
\usetikzlibrary{backgrounds}
\usepackage[letterpaper, margin=1in]{geometry}
\usepackage{multicol}
\usepackage{fourier}  %don't use txfonts b/c the sf doesn't look good

\usepackage{enumitem}
\setlist{topsep=-\parskip/2}


\setlength{\parindent}{0pt}
\setlength{\parskip}{\baselineskip}


\usepackage{mdframed} % makes better frames around examples (eg pagebreak allowed)
  % from mdframed package. used for examples
\newmdenv[
  skipabove=10pt,%\topsep,
  skipbelow=10pt,%\topsep
  ]{allrules}

\usepackage{tcolorbox} % for colored boxes


\usepackage{subfigure}
\usepackage{hyperref}
\hypersetup{colorlinks, urlcolor=red, linkcolor=.} %color urls, but not interal refs
\usepackage{graphicx, floatflt}
\graphicspath{{./figs/}{../figs/}}
\DeclareGraphicsExtensions{.png,.jpg,.pdf}

\usepackage{wrapfig}  %useful for wrapping text around pictures or tables
		      %though it doesn't do a great job inside my example boxes


\definecolor{defbod}{rgb}{.8,.8,.8}     \definecolor{defhead}{rgb}{.7,.1,.1}
\definecolor{thmbod}{rgb}{.8,.8,.8}     \definecolor{thmhead}{rgb}{.1,.1,.9}
\definecolor{factbod}{rgb}{1,1,1}    \definecolor{facthead}{rgb}{.1,.7,.1}
\definecolor{exbod}{rgb}{1,1,1}         \definecolor{exhead}{rgb}{0,.3,0}
\definecolor{soltext}{rgb}{.2,.2,.5}
\definecolor{notebod}{rgb}{.9,.9,.9}    \definecolor{notehead}{rgb}{.1,.1,.7}

\newcommand{\defn}[1]{
 \noindent \colorbox{defbod}{\begin{minipage}{.95\linewidth}
 \hspace{-12pt}\colorbox{defhead}{\sf\color{white} Definition}
 #1
 \end{minipage} } \medskip
 }
\newcommand{\dt}[1]{\textcolor{defhead}{\bfseries #1}}

\newcommand{\thrm}[1]{
 \noindent \colorbox{thmbod}{\begin{minipage}{.95\linewidth}
 \hspace{-12pt}\colorbox{thmhead}{\sf\color{white}Theorem}
 #1
 \end{minipage} } \medskip
 }
\newcommand{\coro}[1]{
 \noindent \colorbox{thmbod}{\begin{minipage}{.95\linewidth}
 \hspace{-12pt}\colorbox{thmhead}{\sf\color{white}Corollary}
 #1
 \end{minipage} } \medskip
 }
\newcommand{\fact}[1]{  %%%% !! note this different from calcnotes
 \noindent \begin{allrules}
 #1
 \end{allrules} \medskip
 }

\newcommand{\note}[1]{
	\noindent \colorbox{notebod}{\begin{minipage}{.95\textwidth}
			\hspace{-12pt}\colorbox{notehead}{\sf\color{white}NOTE}
			#1
	\end{minipage} } \medskip
}


% making the example thing with a line down the left.
\newmdenv[
rightline=false,
bottomline=false,
skipabove=0pt,%\topsep,
skipbelow=0pt,%\topsep
]{siderules}

\newcommand{\ex}[1]{
  \noindent \begin{siderules}
    \textcolor{exhead}{\sf Example}
    #1
  \end{siderules} 
  \medskip
}




%%%%%%%%%% set it up so that these are not displayed when 
%%%%%%%%%% printing notes for students?  or for overhead?
\newcommand{\exsol}[1]{
  {\color{soltext} #1}
}

\newcommand{\mynonumsect}[2][]{
 \cleardoublepage
 \section*{#1\hskip16pt #2}
 \addcontentsline{toc}{section}{#1 #2}
}

\newcommand{\myhead}[1]{
\par%\vskip\baselineskip
\noindent
{\large\bf #1}
}

%%%%%%%%%%%%%%%%%%%%%%%%
% symbol shortcuts

\newcommand{\nice}{\displaystyle}
\newcommand{\eqq}{\ensuremath \stackrel{?}{=}}
\newcommand{\R}{\ensuremath{\mathbb{R}}}




% %%%%%%%%%%%%%%%%%%%%%%%%
% %
% % Creating handouts vs notes
% %% If you use the line below before loading this file then it will be set up for
% %% lecture notes.  Default is for handouts
% %%
% %% \newcommand{\lectureNoteMode}{}
% %%
% %% do not uncomment here, just add the line to the file where you input this one.
% \ifthenelse{\not\isundefined\lectureNoteMode}
% { % lecture notes 
%   \newcommand{\Hvskip}[1]{}
%   \newcommand{\Hvfill}{}
%   \newcommand{\Hpbreak}{}  % pagebreaks for handouts
%   \newcommand{\Npbreak}{\vfill\pagebreak} % pagebreaks for printable notes
%   \newcommand{\Nonly}[1]{#1} % shows in notes, not in handouts
%   \newcommand{\mycleardoublepage}{\cleardoublepage}

%   \newcommand{\ex}[1]{
%    \noindent \begin{allrules}
%     \textcolor{exhead}{\sf Example}\par
%     #1
%     \end{allrules} \medskip
%   }

% } 
% { 
% %% Handout mode 
% %% this section sets it up with smaller margins, and blank spots for writing solns to examples
%   %%%%%%%%%%%%%%%%%%%%%%%%%%%%%%%%%%%%%%%%%%
%   %change frame around examples
%   % from mdframed package
%   \newmdenv[
%   rightline=false,
%   bottomline=false,
%   skipabove=0pt,%\topsep,
%   skipbelow=0pt,%\topsep
%   ]{siderules}
  
%   \newcommand{\ex}[1]{
%     \noindent \begin{siderules}
%       \textcolor{exhead}{\sf Example}
%       #1
%     \end{siderules} 
%     \medskip
%   }
  
%   %
%   %%%%%%%%%%%%%%%%%%%%%%%%%%%%%%%%%%%%%%%%%%
  
%   % don't show solutions for examples
%   \renewcommand{\exsol}[1]{}
  
%   %used for making blank space in handouts. eg: \vs2
%   \newcommand{\Hvskip}[1]{\vskip#1in\null}
%   \newcommand{\Hvfill}{\vfill} %vfill only in handouts mode
%   \newcommand{\Hpbreak}{\pagebreak} % pagebreaks for handouts
%   \newcommand{\Npbreak}{} % pagebreaks for printable notes
%   \newcommand{\Nonly}[1]{} % shows in printable, not in handouts
%   \newcommand{\mycleardoublepage}{\newpage} % don't need blank pages if not printing
  
%   % don't really need margins since not printing.
%   \geometry{margin=0.75in}

%   \renewcommand{\mynonumsect}[2][]{
%     \mycleardoublepage
%     \section*{#1\hskip16pt Handout: #2}
%     \addcontentsline{toc}{section}{#1 #2}
%   }

% } 





\setlength{\parindent}{0pt}

\setcounter{chapter}{1}
\begin{document}

\mynonumsect[1.6]{Modeling with Linear Functions}



\textbf{Goal:} Find linear equations representing real world situations and use them to make predictions.


Model each scenario below with a linear equation.  Then solve the equation and answer the question.

\begin{enumerate}

	\item
	  \emph{Depreciation}. A new car worth \$45,000 is depreciating in value by \$5000 per year.
  
		\begin{enumerate}
	  		\item Find the linear model for the current value of the car, \emph{v}, and the number of years, \emph{y}, after it was purchased.
	  		
	  		\item Interpret the slope and $y$-intercept of the model.
	
	  	
	  		\item If the car is 3 years old, what does the model predict for its value?
	
	
			\item After how many years will the car be worth nothing?
	
	  	\end{enumerate}
 \vfill 
  
  



\pagebreak



	\item Alice is saving money to buy a new smartphone. She currently has \$50 saved and plans to save an additional \$20 each week.
		\begin{enumerate}
			\item Write a linear model for the total amount of money, $M$, she will have saved after $w$ weeks.


			\item Interpret the slope and $y$-intercept of the model.

	
			\item When will she have saved \$300?

		\end{enumerate}
 \vfill


\pagebreak

	\item The percentage of male cigarette smokers in the United States declined from 25.7\% in 2000 \\
	to 16.7\% in 2015. 
		\begin{enumerate}
			\item Find a linear equation relating percentage of male smokers (m) to years since 2000 (t).
			\vfill


			\item Use this equation to predict the year in which the percentage of male smokers \\
			falls below 7\%.
			\vfill


			\item Can you use this equation to estimate the percentage of smokers in the year 2045? Why or why not? 
			\vfill

		\end{enumerate}

\end{enumerate}
  
  
  
  
  
\pagebreak

\defn{\dt{Cost Analysis} \\
	
In a business there are two types of costs: fixed costs and variable costs.  

\begin{itemize}
	\item {\bf Fixed Costs} are costs that do not change with the number of items produced.  \\
	
	Examples of fixed costs: \hrulefill \\
	
	\item {\bf Variable Costs} are costs that change with the number of items produced.  \\
	
	Examples of variable costs: \hrulefill \\
	
	\item The {\bf Total Cost} is the sum of the fixed costs and the variable costs. \\
	\phantom{.}
\end{itemize}	
	
}



\begin{enumerate}[resume]

	\item  A local company produces T-shirts and has a fixed weekly cost of \$15,000 and a variable cost of \$1.50 per shirt.  	
	\begin{enumerate}
		\item Find a linear equation for the cost of producing $x$ shirts? 
		\vfill


		\item What is the cost of producing 10,000 shirts?
		\vfill

		
		\item Interpret the $y$-intercept in the context of this problem.
		\vfill 
	\end{enumerate}
	


	
	
\end{enumerate}






\pagebreak

\myhead{Break-even Analysis}

\defn{\dt{Business Applications} \\
		
	\dt{Revenue} is the amount of money a business receives from selling its products.  This is the number of items sold times the price per item. \\
	
	\dt{Profit} is the difference between the revenue and the total cost.
	\begin{itemize}
		\item If Revenue is greater than Total Cost, then the business makes a Profit.  
		\item If Revenue is less than Total Cost, then the business loses money.  
		\item If Revenue is equal to Total Cost (or equivalently when Profit is zero), then the business breaks even.  \\
	\end{itemize}
	
}




\fact{\dt{Summary of Business Applications Equations: } \\
	
	\begin{enumerate}
		\item \dt{Cost Function} 
		$$C(x) = a+bx$$  
		where $a$ is fixed costs and $b$ is variable cost per item. 
		
		\item \dt{Price-Demand Function} 
		$$p(x) = m-nx$$
		where $x$ is the number of items that can be sold at a price of $p$ dollars per item.
		
		\item \dt{Revenue Function} 
		$$R(x) = xp(x)$$
		where $x$ is the number of items sold at price $p$. 
		
		\item \dt{Profit Function}  
		$$P(x) = R(x) - C(x) = x(m-nx) - (a+bx)$$
		
		\item \dt{Break-Even point} 
		$$C(x)=R(x) \text{\hskip1in or \hskip1in} P(x)=0$$ 
		$$x(m-nx) = (a+bx) \text{\hskip0.5in or \hskip0.5in} x(m-nx) - (a+bx) =0$$
	\end{enumerate}
	\vskip12pt
	Note:  The price-demand funciton is a lower case $p$ and profit is capital $P$.}





\pagebreak


\begin{enumerate}[resume]
	\item A publisher for a promising new novel figures fixed costs (overhead, advances, promotion, copyediting, typesetting, and so on) at \$87,000 and variable costs (printing, paper, binding, shipping) at \$4.50 for each book produced. 
	\begin{enumerate}
  		\item  What is the total cost for $x$ books? \vfill
  	
  		\item If the book is sold to distributors for \$28 each, how many must be produced and sold for the publisher to break even? \vfill  	
	\end{enumerate}
  
	\item A band is planning a concert. The band will receive \$2000 from the venue and \$2 for each ticket sold. The band's expenses are \$2500. How many tickets must be sold for the band to break even? \vfill
	  
	  




\pagebreak



	\item A company manufactures smartphones. Its marketing research department, using statistical techniques, collected the data shown in the table below, where $p$ is the wholesale price per smartphone at which $x$ thousand smartphones can be sold. Using special analytical techniques (regression analysis), an analyst produced the following price-demand function to model the data:
	
	$$p(x)=750-9x \textrm{ for } 0 \leq x \leq 45$$
	

	\begin{enumerate}

		\item What would be the estimated price per smartphone for a demand of 23,000 smartphones?
		\vfill


		\item Write the company's Revenue function and indicate its domain. \vfill


		\item What would be the estimated Revenue for selling 23,000 smartphones?
		\vfill

The financial department for this company established the following cost function for producing and selling $x$ thousand smartphones:
			$$C(x) = 1340+12x$$
		

		\item Use this Cost function and the above Revenue function to write a Profit function for producing and selling $x$ thousand smartphones. Indicate its domain. \vfill


		\item What would be the estimated Profit for selling 23,000 smartphones? \vfill
	\end{enumerate}
	
	
	
\end{enumerate}
	


	


\pagebreak



\myhead{Additional Practice}
\begin{enumerate}
    
  	\item \emph{Health Club Membership}. A health club offers membership for a fee of \$59 plus a monthly fee of \$15 per month.
  		\begin{enumerate}
  			\item Find the linear model for the membership fee, $f$, and the  number of months, $m$, since you have been a member.


		  	\item Interpret the slope of the model.


		  	\item If you have been a member for 24 months, what does the model predict for the fee you have paid so far?

  \end{enumerate}
  \vfill 
 
 
 
 \pagebreak
 
 
 
	\item How many bike computers would a company have to make and sell to break even if the fixed costs are \$36,000, variable costs are \$10.40 per computer, and the computers are sold to retailers for \$15.20 each? \vfill




	\item A company that manufactures and sells x-ray machines has fixed costs of \$1,200,000 and variable costs of \$1,800 per machine. If the machines are sold to hospitals for \$2,400 each, how many machines must be sold for the company to break even? \vfill

\end{enumerate}























\end{document}
